\chapter{CONCLUSION}
In conclusion, the design and development of the automated nutrient dosing and monitoring system of hydroponic crop production demonstrate the sound integration of embedding systems, sensing technologies, and IoT based data management in order to mitigate the drawbacks of traditional hydroponic activities. The developed system has managed to achieve a closed-loop control system to control key nutrient parameters, namely pH and electrical conductivity (EC), which directly affect the availability of nutrients and the dynamics of the normal growth of plants.\\[10pt]
The utilization of an ESP32-WROOM-32 microcontroller as the central processing unit allows efficient handling of real-time sensor data acquisition, control logic execution, and wireless communication. The addition of analog sensors to measure pH and TDS and a digital DHT22 sensor for environmental monitoring ensures the complete monitoring of the nature of the nutrient solution and surrounding environment. Furthermore, the implementation of temperature compensation in TDS and pH calculations also improve the accuracy of measurements, including the intrinsic dependence and environmental changes.\\[10pt]
An important advancement made by this system is the implementation of a threshold-based closed-loop control strategy with stabilization delay, which is especially most appropriate to use in the hydroponic systems with a slow dynamic response. Compared to the traditional PID-based techniques, the offered method is simpler to compute and avoids issues such as overshoot and oscillations, making it highly suitable for low-cost microcontroller platforms. The decision to regulate the EC first as opposed to pH adjustment is also an indication of the view of the relationship between the concentration of the nutrient and the acidity of the solution hence enhancing stability and effectiveness.\\[10pt]
The integration of IoT capabilities through Wi-Fi communication and cloud-based database storage enables real-time data logging, remote monitor, and transparency of the system. The backend framework is based on Django and RESTful API to ensure safe data exchange and scalability of the system, whereas the front end interface ensures the intuitive display of sensor measurements and system conditions. This end to end architecture confirms the possibility of implementing smart hydroponic systems in real life environments with little human intervention.\\[10pt]
Experimental observations confirm that the system is capable of continuously sensing, transmitting, storing, and displaying critical parameters, thereby establishing a stable monitoring structure. Though some hardware limitations like the instability of ADC pin and sensor drift during testing were found, these issues can be used in the optimization of the system in the future, like a better hardware design and better calibration processes.\\[10pt]
Despite the limitation of being a single crop operation and manually calibrated, the system has great potential of being extended further in the future. Multi crop profiles, adaptive control strategies, and machine learning-based predictive models might help change the system to a proactive rather than a reactive control framework. Also, more sensor redundancy, fault detection and mobile-based user interfaces can be enhanced to improve the system reliability and usability.\\[10pt]
Overall, this project demonstrates a scalable, cost-effective, and technically feasible project to utilize precision agriculture. Through automated systems, real-time monitoring, and smart control, it helps to enhance sustainable hydroponic farming practices, especially in resource- constrained and urban environments.
